\chapter{Modelica Concrete Syntax}\label{modelica-concrete-syntax}

\section{Lexical conventions}\label{lexical-conventions}

The following syntactic metasymbols are used (extended BNF):
\begin{center}
\begin{tabular}{c l}
\hline
\tablehead{Syntax} & \tablehead{Description}\\
\hline
\hline
{\lstinline[language=grammar]![ $\ldots$ ]!} & Optional\\
{\lstinline[language=grammar]!{ $\ldots$ }!} & Repeat zero or more times\\
{\lstinline[language=grammar]!$\ldots$ | $\ldots$!} & Alternatives\\
{\lstinline[language=grammar]!"$\mathit{text}$"!} & The $\mathit{text}$ is treated as a single token (no white-space between any characters)\\
\hline
\end{tabular}
\end{center}

The following lexical units are defined:
% Beware that the first lines of the lexing rules below are duplicated in \cref{identifiers}, and must be kept in sync.

\begin{lexicalunit}{IDENT}
\begin{lstlisting}
NON-DIGIT { DIGIT | NON-DIGIT } | Q-IDENT
\end{lstlisting}
\end{lexicalunit}

\begin{lexicalunit}{Q-IDENT}
\begin{lstlisting}
"'" { Q-CHAR | S-ESCAPE } "'"
\end{lstlisting}
\end{lexicalunit}

\begin{lexicalunit}{NON-DIGIT}
\begin{lstlisting}
"_" | letters "a" $\ldots$ "z" | letters "A" $\ldots$ "Z"
\end{lstlisting}
\end{lexicalunit}

\begin{lexicalunit}{DIGIT}
\begin{lstlisting}
"0" | "1" | "2" | "3" | "4" | "5" | "6" | "7" | "8" | "9"
\end{lstlisting}
\end{lexicalunit}

\begin{lexicalunit}{Q-CHAR}
\begin{lstlisting}
NON-DIGIT | DIGIT
| " " | """ | "." | "," | ":" | ";" | "!" | "?" | "@" | "#" | "$\mbox{\textdollar}$"
| "+" | "-" | "*" | "/" | "^" | "=" | "&" | "|" | "~" | "%"
| "<" | ">" | "(" | ")" | "[" | "]" | "{" | "}"
\end{lstlisting}
\end{lexicalunit}

\begin{lexicalunit}{S-ESCAPE}
\begin{lstlisting}
"\'" | "\"" | "\?" | "\\"
| "\a" | "\b" | "\f" | "\n" | "\r" | "\t" | "\v"
\end{lstlisting}
\end{lexicalunit}

\begin{lexicalunit}{STRING}
\begin{lstlisting}
""" { S-CHAR | S-ESCAPE } """
\end{lstlisting}
\end{lexicalunit}

\begin{lexicalunit}{S-CHAR}
\begin{lstlisting}
see below
\end{lstlisting}
\end{lexicalunit}

\begin{lexicalunit}{UNSIGNED-INTEGER}
\begin{lstlisting}
DIGIT { DIGIT }
\end{lstlisting}
\end{lexicalunit}

\begin{lexicalunit}{UNSIGNED-REAL}
\begin{lstlisting}
UNSIGNED-INTEGER  "." [ UNSIGNED-INTEGER ]
| UNSIGNED_INTEGER [ "." [ UNSIGNED_INTEGER ] ]
  ( "e" | "E" ) [ "+" | "-" ] UNSIGNED-INTEGER
| "."  UNSIGNED-INTEGER [ ( "e" | "E" ) [ "+" | "-" ] UNSIGNED-INTEGER ]
\end{lstlisting}
\end{lexicalunit}

\lexicalunitref{S-CHAR} is any member of the Unicode character set (\url{https://unicode.org}; see \cref{mapping-package-class-structures-to-a-hierarchical-file-system} for storing as UTF-8 on files) except double-quote `"', and backslash `\textbackslash{}'.

For identifiers the redundant escapes (`\lstinline!\?!' and `\lstinline!\"!') are the same as the corresponding non-escaped variants (`\lstinline!?!' and '\lstinline!"!').
The single quotes are part of an identifier.
For example, the identifiers \lstinline!'x'! and \lstinline!x! are different.

Note:
\begin{itemize}
\item
  White-space and comments can be used between separate lexical units and/or symbols, and also separates them.
  Each lexical unit will consume the maximum number of characters from the input stream.
  White-space and comments cannot be used inside other lexical units, except for \lexicalunitref{STRING} and \lexicalunitref{Q-IDENT} where they are treated as part of the \lexicalunitref{STRING} or \lexicalunitref{Q-IDENT} lexical unit.
\item
  Concatenation of string literals requires a binary expression.
  For example, \lstinline!"a" + "b"! evaluates to \lstinline!"ab"!.
  There is no support for the C/C++ style of concatenating adjacent string literal tokens (for example, \lstinline[language=C]!"a" "b"! becoming \lstinline[language=C]!"ab"!).
\item
  Modelica uses the same comment syntax as C++ and Java (i.e., \lstinline!//! signals the start of a line comment and \lstinline!/* $\ldots$ */! is a multi-line comment); comments may contain any Unicode character.
  Modelica also has structured comments in the form of annotations and string comments.
\item
  In the grammar, keywords of the Modelica language are highlighted with color, for example, \lstinline[language=grammar]!equation!.
\item
  Productions use hyphen as separator both in the grammar and in the text.
  (Previously the grammar used underscore.)
\end{itemize}

If a description-string starts with the tag \lstinline!<html>! or \lstinline!<HTML>!, the entire string is HTML encoded (assumed to end with \lstinline!</html>! or \lstinline!</HTML>!, but to be rendered as HTML even if the end-tag is missing).
Otherwise, the entire string is rendered as is.
HTML encoded content may contain links in the same way as class documentation, see \cref{documentation}.


\section{Grammar}\label{grammar}

\subsection{Stored Definition -- Within}\label{stored-definition-within}

\begin{production}{stored-definition}
\begin{lstlisting}
[ within-clause ]
{ [ final ] class-definition ";" }
\end{lstlisting}
\end{production}

\begin{production}{within-clause}
\begin{lstlisting}
within [ name ] ";"
\end{lstlisting}
\end{production}


\subsection{Class Definition}\label{class-definition}

\begin{production}{class-definition}
\begin{lstlisting}
[ encapsulated ] class-prefixes class-specifier
\end{lstlisting}
\end{production}

\begin{production}{class-prefixes}
\begin{lstlisting}
[ partial ]
( class
  | model
  | [ operator ] record
  | block
  | [ expandable ] connector
  | type
  | package
  | [ pure | impure ] [ operator ] function
  | operator
)
\end{lstlisting}
\end{production}

\begin{production}{class-specifier}
\begin{lstlisting}
long-class-specifier | short-class-specifier | der-class-specifier
\end{lstlisting}
\end{production}

\begin{production}{long-class-specifier}
\begin{lstlisting}
IDENT description-string composition end IDENT
| extends IDENT [ class-modification ] description-string composition
  end IDENT
\end{lstlisting}
\end{production}

\begin{production}{short-class-specifier}
\begin{lstlisting}
IDENT "=" base-prefix type-specifier [ array-subscripts ]
[ class-modification ] description
| IDENT "=" enumeration "(" ( [ enum-list ] | ":" ) ")" description
\end{lstlisting}
\end{production}

\begin{production}{der-class-specifier}
\begin{lstlisting}
IDENT "=" der "(" type-specifier "," IDENT { "," IDENT } ")" description
\end{lstlisting}
\end{production}

\begin{production}{base-prefix}
\begin{lstlisting}
[ input | output ]
\end{lstlisting}
\end{production}

\begin{production}{enum-list}
\begin{lstlisting}
enumeration-literal { "," enumeration-literal }
\end{lstlisting}
\end{production}

\begin{production}{enumeration-literal}
\begin{lstlisting}
IDENT description
\end{lstlisting}
\end{production}

\begin{production}{composition}
\begin{lstlisting}
element-list
{ public element-list
  | protected element-list
  | equation-section
  | algorithm-section
}
[ external [ language-specification ]
  [ external-function-call ] [ annotation-clause ] ";"
]
[ annotation-clause ";" ]
\end{lstlisting}
\end{production}

\begin{production}{language-specification}
\begin{lstlisting}
STRING
\end{lstlisting}
\end{production}

\begin{production}{external-function-call}
\begin{lstlisting}
[ component-reference "=" ]
IDENT "(" [ expression-list ] ")"
\end{lstlisting}
\end{production}

\begin{production}{element-list}
\begin{lstlisting}
{ element ";" }
\end{lstlisting}
\end{production}

\begin{production}{element}
\begin{lstlisting}
import-clause
| extends-clause
| [ redeclare ]
  [ final ]
  [ inner ] [ outer ]
  ( class-definition
    | component-clause
    | replaceable ( class-definition | component-clause )
      [ constraining-clause description ]
  )
\end{lstlisting}
\end{production}

\begin{production}{import-clause}
\begin{lstlisting}
import
( IDENT "=" name
  | name [ ".*" | "." ( "*" | "{" import-list "}" ) ]
)
description
\end{lstlisting}
\end{production}

\begin{production}{import-list}
\begin{lstlisting}
IDENT { "," IDENT }
\end{lstlisting}
\end{production}


\subsection{Extends}\label{extends}

\begin{production}{extends-clause}
\begin{lstlisting}
extends type-specifier [ class-or-inheritance-modification ] [ annotation-clause ]
\end{lstlisting}
\end{production}

\begin{production}{constraining-clause}
\begin{lstlisting}
constrainedby type-specifier [ class-modification ]
\end{lstlisting}
\end{production}

\begin{production}{class-or-inheritance-modification}
\begin{lstlisting}
"(" [ argument-or-inheritance-modification-list ] ")"
\end{lstlisting}
\end{production}

\begin{production}{argument-or-inheritance-modification-list}
\begin{lstlisting}
 ( argument | inheritance-modification ) { "," ( argument | inheritance-modification ) }
\end{lstlisting}
\end{production}

\begin{production}{inheritance-modification}
\begin{lstlisting}
 break ( connect-equation | IDENT )
\end{lstlisting}
\end{production}


\subsection{Component Clause}\label{component-clause}

\begin{production}{component-clause}
\begin{lstlisting}
type-prefix type-specifier [ array-subscripts ] component-list
\end{lstlisting}
\end{production}

\begin{production}{type-prefix}
\begin{lstlisting}
[ flow | stream ]
[ discrete | parameter | constant ]
[ input | output ]
\end{lstlisting}
\end{production}

\begin{production}{component-list}
\begin{lstlisting}
component-declaration { "," component-declaration }
\end{lstlisting}
\end{production}

\begin{production}{component-declaration}
\begin{lstlisting}
declaration [ condition-attribute ] description
\end{lstlisting}
\end{production}

\begin{production}{condition-attribute}
\begin{lstlisting}
if expression
\end{lstlisting}
\end{production}

\begin{production}{declaration}
\begin{lstlisting}
IDENT [ array-subscripts ] [ modification ]
\end{lstlisting}
\end{production}


\subsection{Modification}\label{modification}

\begin{production}{modification}
\begin{lstlisting}
class-modification [ "=" modification-expression ]
| "=" modification-expression
\end{lstlisting}
\end{production}

\begin{production}{modification-expression}
\begin{lstlisting}
expression
| break
\end{lstlisting}
\end{production}

\begin{production}{class-modification}
\begin{lstlisting}
"(" [ argument-list ] ")"
\end{lstlisting}
\end{production}

\begin{production}{argument-list}
\begin{lstlisting}
argument { "," argument }
\end{lstlisting}
\end{production}

\begin{production}{argument}
\begin{lstlisting}
element-modification-or-replaceable
| element-redeclaration
\end{lstlisting}
\end{production}

\begin{production}{element-modification-or-replaceable}
\begin{lstlisting}
[ each ] [ final ] ( element-modification | element-replaceable )
\end{lstlisting}
\end{production}

\begin{production}{element-modification}
\begin{lstlisting}
name [ modification ] description-string
\end{lstlisting}
\end{production}

\begin{production}{element-redeclaration}
\begin{lstlisting}
redeclare [ each ] [ final ]
( short-class-definition | component-clause1 | element-replaceable )
\end{lstlisting}
\end{production}

\begin{production}{element-replaceable}
\begin{lstlisting}
replaceable ( short-class-definition | component-clause1 )
[ constraining-clause ]
\end{lstlisting}
\end{production}

\begin{production}{component-clause1}
\begin{lstlisting}
type-prefix type-specifier component-declaration1
\end{lstlisting}
\end{production}

\begin{production}{component-declaration1}
\begin{lstlisting}
declaration description
\end{lstlisting}
\end{production}

\begin{production}{short-class-definition}
\begin{lstlisting}
class-prefixes short-class-specifier
\end{lstlisting}
\end{production}


\subsection{Equations}\label{equations1}

\begin{production}{equation-section}
\begin{lstlisting}
[ initial ] equation { some-equation ";" }
\end{lstlisting}
\end{production}

\begin{production}{algorithm-section}
\begin{lstlisting}
[ initial ] algorithm { statement ";" }
\end{lstlisting}
\end{production}

\begin{production}{some-equation}
\begin{lstlisting}
( equation-or-procedure
  | if-equation
  | for-equation
  | connect-equation
  | when-equation
)
description
\end{lstlisting}
\end{production}

\begin{production}{equation-or-procedure}
\begin{lstlisting}
simple-expression "=" expression
| function-call
\end{lstlisting}
\end{production}

\begin{production}{statement}
\begin{lstlisting}
( statement-or-procedure
  | "(" output-expression-list ")" ":="
    function-call
  | break
  | return
  | if-statement
  | for-statement
  | while-statement
  | when-statement
)
description
\end{lstlisting}
\end{production}

\begin{production}{statement-or-procedure}
\begin{lstlisting}
component-reference ":=" expression
| function-call
\end{lstlisting}
\end{production}

\begin{production}{function-call}
\begin{lstlisting}
component-reference function-call-args
\end{lstlisting}
\end{production}

\begin{production}{if-equation}
\begin{lstlisting}
if expression then
  { some-equation ";" }
{ elseif expression then
  { some-equation ";" }
}
[ else
  { some-equation ";" }
]
end if
\end{lstlisting}
\end{production}

\begin{production}{if-statement}
\begin{lstlisting}
if expression then
  { statement ";" }
{ elseif expression then
  { statement ";" }
}
[ else
  { statement ";" }
]
end if
\end{lstlisting}
\end{production}

\begin{production}{for-equation}
\begin{lstlisting}
for for-indices loop
  { some-equation ";" }
end for
\end{lstlisting}
\end{production}

\begin{production}{for-statement}
\begin{lstlisting}
for for-indices loop
  { statement ";" }
end for
\end{lstlisting}
\end{production}

\begin{production}{for-indices}
\begin{lstlisting}
for-index { "," for-index }
\end{lstlisting}
\end{production}

\begin{production}{for-index}
\begin{lstlisting}
IDENT [ in expression ]
\end{lstlisting}
\end{production}

\begin{production}{while-statement}
\begin{lstlisting}
while expression loop
  { statement ";" }
end while
\end{lstlisting}
\end{production}

\begin{production}{when-equation}
\begin{lstlisting}
when expression then
  { some-equation ";" }
{ elsewhen expression then
  { some-equation ";" }
}
end when
\end{lstlisting}
\end{production}

\begin{production}{when-statement}
\begin{lstlisting}
when expression then
  { statement ";" }
{ elsewhen expression then
  { statement ";" }
}
end when
\end{lstlisting}
\end{production}

\begin{production}{connect-equation}
\begin{lstlisting}
connect "(" component-reference "," component-reference ")"
\end{lstlisting}
\end{production}

\begin{nonnormative}
The productions \productionref{equation-or-procedure} and \productionref{statement-or-procedure} are not suitable for recursive descent parsers.
A work-around is to left-factor them and for \productionref{equation-or-procedure} introduce semantic checks to ensure that only the grammar above is accepted.
\end{nonnormative}


\subsection{Expressions}\label{expressions1}

\begin{production}{expression}
\begin{lstlisting}
simple-expression
| if expression then expression
  { elseif expression then expression }
  else expression
\end{lstlisting}
\end{production}

\begin{production}{simple-expression}
\begin{lstlisting}
logical-expression [ ":" logical-expression [ ":" logical-expression ] ]
\end{lstlisting}
\end{production}

\begin{production}{logical-expression}
\begin{lstlisting}
logical-term { or logical-term }
\end{lstlisting}
\end{production}

\begin{production}{logical-term}
\begin{lstlisting}
logical-factor { and logical-factor }
\end{lstlisting}
\end{production}

\begin{production}{logical-factor}
\begin{lstlisting}
[ not ] relation
\end{lstlisting}
\end{production}

\begin{production}{relation}
\begin{lstlisting}
arithmetic-expression [ relational-operator arithmetic-expression ]
\end{lstlisting}
\end{production}

\begin{production}{relational-operator}
\begin{lstlisting}
"<" | "<=" | ">" | ">=" | "==" | "<>"
\end{lstlisting}
\end{production}

\begin{production}{arithmetic-expression}
\begin{lstlisting}
[ add-operator ] term { add-operator term }
\end{lstlisting}
\end{production}

\begin{production}{add-operator}
\begin{lstlisting}
"+" | "-" | ".+" | ".-"
\end{lstlisting}
\end{production}

\begin{production}{term}
\begin{lstlisting}
factor { mul-operator factor }
\end{lstlisting}
\end{production}

\begin{production}{mul-operator}
\begin{lstlisting}
"*" | "/" | ".*" | "./"
\end{lstlisting}
\end{production}

\begin{production}{factor}
\begin{lstlisting}
primary [ ( "^" | ".^" ) primary ]
\end{lstlisting}
\end{production}

\begin{production}{primary}
\begin{lstlisting}
unsigned-number
| STRING
| false
| true
| time
| ( component-reference | der | initial | pure ) function-call-args
| component-reference
| "(" output-expression-list ")" [ ( array-subscripts | "." IDENT ) ]
| "[" expression-list { ";" expression-list } "]"
| "{" array-arguments "}"
| end
\end{lstlisting}
\end{production}

\begin{production}{unsigned-number}
\begin{lstlisting}
UNSIGNED-INTEGER | UNSIGNED-REAL
\end{lstlisting}
\end{production}

\begin{production}{type-specifier}
\begin{lstlisting}
["."] name
\end{lstlisting}
\end{production}

\begin{production}{name}
\begin{lstlisting}
IDENT { "." IDENT }
\end{lstlisting}
\end{production}

\begin{production}{component-reference}
\begin{lstlisting}
[ "." ] IDENT [ array-subscripts ] { "." IDENT [ array-subscripts ] }
\end{lstlisting}
\end{production}

\begin{production}{result-reference}
\begin{lstlisting}
component-reference
| time
| der "(" ( component-reference | time ) [ "," UNSIGNED-INTEGER ] ")"
\end{lstlisting}
\end{production}

\begin{production}{function-call-args}
\begin{lstlisting}
"(" [ function-arguments ] ")"
\end{lstlisting}
\end{production}

\begin{production}{function-arguments}
\begin{lstlisting}
expression [ "," function-arguments-non-first | for for-indices ]
| function-partial-application [ "," function-arguments-non-first ]
| named-arguments
\end{lstlisting}
\end{production}

\begin{production}{function-arguments-non-first}
\begin{lstlisting}
function-argument [ "," function-arguments-non-first ]
| named-arguments
\end{lstlisting}
\end{production}

\begin{production}{array-arguments}
\begin{lstlisting}
expression [ "," array-arguments-non-first | for for-indices ]
\end{lstlisting}
\end{production}

\begin{production}{array-arguments-non-first}
\begin{lstlisting}
expression [ "," array-arguments-non-first ]
\end{lstlisting}
\end{production}

\begin{production}{named-arguments}
\begin{lstlisting}
named-argument [ "," named-arguments ]
\end{lstlisting}
\end{production}

\begin{production}{named-argument}
\begin{lstlisting}
IDENT "=" function-argument
\end{lstlisting}
\end{production}

\begin{production}{function-argument}
\begin{lstlisting}
function-partial-application | expression
\end{lstlisting}
\end{production}

\begin{production}{function-partial-application}
\begin{lstlisting}
function type-specifier "(" [ named-arguments ] ")"
\end{lstlisting}
\end{production}

\begin{production}{output-expression-list}
\begin{lstlisting}
[ expression ] { "," [ expression ] }
\end{lstlisting}
\end{production}

\begin{production}{expression-list}
\begin{lstlisting}
expression { "," expression }
\end{lstlisting}
\end{production}

\begin{production}{array-subscripts}
\begin{lstlisting}
"[" subscript { "," subscript } "]"
\end{lstlisting}
\end{production}

\begin{production}{subscript}
\begin{lstlisting}
":" | expression
\end{lstlisting}
\end{production}

\begin{production}{description}
\begin{lstlisting}
description-string [ annotation-clause ]
\end{lstlisting}
\end{production}

\begin{production}{description-string}
\begin{lstlisting}
[ STRING { "+" STRING } ]
\end{lstlisting}
\end{production}

\begin{production}{annotation-clause}
\begin{lstlisting}
annotation class-modification
\end{lstlisting}
\end{production}
